El presente anexo detalla los pasos necesarios para obtener el código fuente de la aplicación LogArt, configurar el entorno y ejecutar la aplicación en un entorno local utilizando Docker y Docker Compose.

\section{Requisitos Previos}
\label{appA:requisitosPrevios}

Antes de continuar, asegúrese de tener instaladas las siguientes herramientas en el sistema:
\begin{itemize}
    \item \textbf{Git:} Para clonar el repositorio del proyecto. (Guía de instalación disponible en \url{https://git-scm.com/book/en/v2/Getting-Started-Installing-Git})
    \item \textbf{Docker Engine:} Para construir y ejecutar los contenedores de la aplicación. (Guía de instalación: \url{https://docs.docker.com/get-docker/})
    \item \textbf{Docker Compose:} Para orquestar la aplicación multi-contenedor. (Guía de instalación: \url{https://docs.docker.com/compose/install/})
\end{itemize}
Se recomienda utilizar versiones recientes de Docker y Docker Compose (ej. Docker +v24.x, Docker Compose con sintaxis +v3.8).

\section{Obtención del Código Fuente}
\label{appA:obtencionCodigo}

El código fuente de LogArt está alojado en un repositorio de GitHub. Para obtenerlo, clone el repositorio utilizando el siguiente comando en su terminal Git:
\begin{verbatim}
    git clone https://github.com/codeurjc-students/2024-logart.git
\end{verbatim}

Una vez clonado, navegue al directorio específico donde se encuentra el archivo \texttt{Dockerfile}, que para este proyecto es \texttt{LogArtApp/docker/} dentro del repositorio clonado:
\begin{verbatim}
    cd 2024-logart/LogArtApp/docker/
\end{verbatim}

\section{Configuración del Entorno (\texttt{.env})}
\label{appA:configuracionEntorno}
LogArt requiere un archivo de configuración de entorno llamado \texttt{.env} para definir variables críticas como las credenciales de la base de datos, configuración de correo electrónico y secretos de la aplicación. Este archivo debe ser creado por el usuario que va a construir e instalar la aplicación, en el directorio \texttt{LogArtApp/docker/} (el mismo directorio donde se encuentra el archivo \texttt{docker-compose.yml}.

Cree entonces un archivo llamado \texttt{.env} en ese directorio y añada el siguiente contenido, \textbf{reemplazando los valores placeholder (your\_...)} con sus propias credenciales o valores deseados:

\begin{verbatim}
# MongoDB Configuration (credenciales para el usuario root de MongoDB)
MONGO_INITDB_ROOT_USERNAME=your_mongo_root_user
MONGO_INITDB_ROOT_PASSWORD=your_mongo_root_password

# Backend Database Connection String (usa las credenciales anteriores)
DB_CONNECTION_STRING=mongodb://your_mongo_root_user
:your_mongo_root_password@mongo:27017/logartdb?authSource=admin

# JWT Configuration (usar cadenas largas y aleatorias)
ACCESS_TOKEN_SECRET=your_very_strong_and_random_access_token_secret

# Email Configuration (para verificación de cuenta y recuperación de contraseña)
EMAIL_USER=your_gmail_address_for_sending_emails@gmail.com
EMAIL_PASS=your_gmail_app_password_or_regular_password 
# Nota: Para Gmail, se debe usar una "Contraseña de aplicación".

# Application Configuration
BASE_URL=https://localhost:8443 # URL base donde la aplicación será accesible
PORT=8443 # Puerto en el que escuchará el servidor Node.js dentro del contenedor
NODE_ENV=production # Entorno de ejecución de Node.js

# SSL Configuration (rutas relativas al directorio del backend dentro del contenedor)
SSL_KEY_PATH=ssl/server.key
SSL_CERT_PATH=ssl/server.cert
\end{verbatim}

\textbf{Nota Importante sobre SSL para Ejecución Local:}
La aplicación está configurada para utilizar HTTPS. El \texttt{Dockerfile} (etapa \texttt{backend-build}) no incluye la generación de certificados SSL autofirmados, estos deben ser creados y proporcionados. Para desarrollo local, deberá:
\begin{enumerate}
    \item Crear un directorio \textbf{ssl} dentro de \texttt{LogArtApp/backend/}.
    \item Generar sus propios archivos \texttt{server.key} y \texttt{server.cert} autofirmados (por ejemplo, usando OpenSSL) y colocarlos en la carpeta mencionada (\texttt{LogArtApp/backend/ssl/}.
\end{enumerate}
Alternativamente, si no se desea usar HTTPS, ni configurar SSL, y únicamente para un entorno de desarrollo local estricto, sería necesario modificar el código del servidor Express.js (\texttt{backend/server.js}) para que escuche sobre HTTP en lugar de HTTPS, lo cual no es recomendado. \textbf{No se asegura el correcto funcionamiento de la web y sus contenidos asociados usando HTTP}.

\section{Ejecución de la Aplicación}
\label{appA:ejecucionApp}

Una vez clonado el repositorio, configurado el archivo \texttt{.env} en el directorio \texttt{LogArtApp/docker/}, y preparados los certificados SSL, existen dos métodos principales para ejecutar la aplicación LogArt utilizando Docker.

\subsection{Método 1: Construcción Manual de la Imagen y Ejecución con Docker Compose}
\label{appA:ejecucionMetodo1}
Este método es útil si se desea construir la imagen Docker localmente a partir del código fuente más reciente o con modificaciones propias antes de levantar los servicios.
\begin{enumerate}
    \item \textbf{Construir la Imagen Docker Localmente:}
    Navegue al directorio raíz del proyecto clonado (\texttt{2024-logart/LogArtApp/}). Desde allí, ejecute el siguiente comando para construir la imagen Docker utilizando el \texttt{Dockerfile} ubicado en la carpeta \texttt{LogArtApp/docker/Dockerfile}. Puede asignar el tag que desee:
    \begin{verbatim}
        docker build -t davidmorenoo/logartapp:latestManual -f docker/Dockerfile .
    \end{verbatim}
    \textbf{Nota:} Si utiliza un tag diferente al que está especificado en el archivo \texttt{docker-compose.yml}, deberá actualizar la directiva \texttt{image} del servicio \texttt{app} en el archivo \texttt{docker-compose.yml} para que coincida con el tag de la imagen que acaba de construir.
    \item \textbf{Levantar los Servicios con Docker Compose:}
    Navegue al directorio \texttt{LogArtApp/docker/} (donde se encuentra el archivo \texttt{docker-compose.yml} y su \texttt{.env} configurado). Ejecute el siguiente comando para iniciar todos los servicios (la aplicación LogArt y la base de datos MongoDB):
    \begin{verbatim}
        docker compose up
    \end{verbatim}
\end{enumerate}

\subsection{Método 2: Ejecución Directa con Docker Compose (Utilizando Imagen Pre-construida}
\label{appA:ejecucionMetodo2}
Este método es el recomendado, ya que es más directo, al estar el \texttt{docker-compose.yml} configurado para utilizar una imagen pre-construida de DockerHub.
\begin{enumerate}
    \item \textbf{Asegurar Configuración del \texttt{docker-compose.yml}:}
    Verifique que el archivo \texttt{docker-compose} esté correctamente configurado. Por defecto, el servicio \texttt{app} utiliza la imagen \texttt{latestManual} disponible en el repositorio público \texttt{davidmorenoo/logartapp} de DockerHub.
    \item \textbf{Levantar los Servicios con Docker Compose:}
    Navegue al directorio \texttt{LogArtApp/docker/} y ejecute el siguiente comando:
    \begin{verbatim}
        docker compose up
    \end{verbatim}
\end{enumerate}

\subsection{Acceso a la Aplicación}
\label{appA:accesoApp}
Independientemente del método de ejecución elegido, una vez que los contenedores estén en funcionamiento, la aplicación LogArt estará disponible en el navegador web a través de la URL:
\begin{center}
    \url{https://localhost:8443}
\end{center}
\textbf{Nota sobre Certificados SSL Autofirmados:} Al utilizar certificados SSL autofirmados para el desarrollo local sobre HTTPS, es probable que su navegador muestre una advertencia de seguridad. Deberá aceptar a través de las opciones de ``configuración avanzada'' para acceder al sitio (generalmente buscando una opción como ``Acceder a localhost (sitio no seguro)'' o similar).

\section{Detener la Aplicación}
\label{appA:detenerApp}
Para detener todos los servicios de la aplicación (LogArt y MongoDB) y eliminar los contenedores, navegue al directorio \texttt{LogArtApp/docker/} y ejecute:
\begin{verbatim}
    docker compose down
\end{verbatim}
Si desea volver a iniciar la aplicación, puede ejecutar nuevamente \texttt{docker compose up}.

\textbf{Credenciales de Acceso de Ejemplo y Datos Precargados}
\label{appA:credencialesDatos}
LogArt se inicia con un conjunto de datos de ejemplo y credenciales predefinidas para facilitar la exploración y testeo inicial:

\begin{itemize}
    \item \textbf{Usuario Administrador:}
    \begin{itemize}
        \item \textbf{Username/Email:}
        \texttt{admin@gmail.com}
        \item \textbf{Password:} \texttt{admin123}
        \item \textbf{Permisos:} Acceso completo a todas las funcionalidades, incluyendo el panel de administración, la moderación de contenido y el resto de funcionalidades definidas para usuario administrador en \ref{ssec:actoresSistema}
    \end{itemize}
    \item \textbf{Usuario Registrado:}
    \begin{itemize}
        \item \textbf{Username/Email:}
        \texttt{pepe@gmail.com}
        \item \textbf{Password:} \texttt{hola123}
        \item \textbf{Permisos:} Creación y gestión de sus propios objetos y comentarios, junto con el resto de funcionalidades definidas para usuario registrado en \ref{ssec:actoresSistema}
    \end{itemize}
    \item \textbf{Usuario No Registrado:}
        \begin{itemize}
            \item Puede visualizar la página principal y los objetos que hayan sido compartidos públicamente por otros usuarios, junto con el resto de funcionalidades definidas para usuario no registrado en \ref{ssec:actoresSistema}
        \end{itemize}
\end{itemize}

La aplicación también incluye datos de ejemplo preconfigurados para las disciplinas (Videojuegos, Libros, Canciones), algunos objetos y comentarios, que pueden ser modificados o eliminados por los usuarios o por el administrador.

\section{Interacción con la Base de Datos (Opcional para Desarrollo Avanzado)}
\label{appA:interaccionBD}
Si desea realizar depuración avanzada, y necesita interactuar con la base de datos MongoDB que se ejecuta en el contenedor, puede hacerlo. Con los contenedores en ejecución (tras un \texttt{docker compose up}, abra una terminal y ejecute el siguiente comando para acceder al shell de MongoDB dentro del contenedor \texttt{mongo}. Deberá reemplazar \texttt{your\_mongo\_root\_user} y \texttt{your\_mongo\_root\_password} con los valores que configuró en su archivo \texttt{.env}:
\begin{verbatim}
    docker exec -it docker-mongo-1 mongosh \
   --username your_mongo_root_user \
   --password your_mongo_root_password \
   --authenticationDatabase admin
\end{verbatim}
Una vez dentro del shell de \texttt{mongosh}, puede ejecutar comandos de MongoDB. Por ejemplo, para ver los objetos en la base de datos de LogArt:
\begin{verbatim}
    use logartdb;
    show collections;
    db.objects.find().pretty();
\end{verbatim}
Este acceso directo a la base de datos debe realizarse con cuidado, y solo en casos de depuración avanzada.